\chapter*{ABSTRACT}
\addcontentsline{toc}{chapter}{Abstract}
\thispagestyle{plain}
\noindent\rule{\linewidth}{0.5pt}

Large Language Models (LLMs) have become capable tools for competitive programming (CP) assistance, yet existing evaluation benchmarks reduce model performance to a single aggregate score. This obscures a more informative question: which categories of algorithmic bugs can a model reliably expose, and where does its reasoning break down? This paper presents AlgoBugs, a diagnostic benchmark that addresses this gap through structured fault-exposure evaluation. AlgoBugs consists of 1,000 paired C++ submissions drawn from 40 Codeforces problems across eight difficulty brackets (rating 1400--2100). Each pair contains a Wrong Answer submission and its corresponding Accepted solution for the same problem, with the bug manually identified and classified under a novel eight-category taxonomy (T1: Integer Overflow, T2: Modular Arithmetic, T3: Off-by-One, T4: Wrong Conditional, T5: Algorithmic Flaw, T6: State/Initialization, T7: Corner Case, T8: I/O Format). Taxonomic labels were validated using inter-rater reliability analysis (Cohen's $\kappa$). Five LLMs spanning frontier, mid-tier, open-source, and code-specialist tiers were evaluated under three prompting strategies: zero-shot, chain-of-thought, and few-shot. Each model was tasked with generating a test input that causes the buggy solution to produce a wrong answer while the correct solution does not. Test case validity and fault exposure were determined by a local execution-based differential testing engine using g++ C++20 compilation. Results show Fault Exposure Rates (FER) ranging from 14\% to 21\% across all models, with T3 (Off-by-One) and T4 (Wrong Conditional) being the most exposed categories and T1 (Integer Overflow) and T2 (Modular Arithmetic) being the hardest. Chain-of-thought prompting provided no consistent improvement over zero-shot. Few-shot prompting improved FER for LLaMA 3.3 70B (+3.4 percentage points) and GPT-5.1 (+5.6 percentage points) but not for the remaining three models, indicating model-specific sensitivity to task demonstrations rather than a general benefit. These findings reveal that current LLMs generate specification-valid test inputs rather than targeted fault-exposing ones, highlighting a gap between aggregate benchmark scores and genuine diagnostic reasoning capability.

